\chapter{Structures algébriques}\label{ch:b1:structures}

Les mêmes règles de calcul ne cessent de réapparaître~: entiers, nombres
réels, nombres complexes, classes de \omterm{def:b1:arith:congruence}{congruence}, et bientôt polynômes
(\cref{ch:b1:poly}), vecteurs et matrices
(\cref{ch:b1:vspaces,ch:b1:matrices}). L'algèbre extrait les schémas
communs et leur donne un nom~: \emph{\omterm{def:b1:structures:group}{groupe}}, \emph{\omterm{def:b1:structures:ring}{anneau}},
\emph{\omterm{def:b1:structures:field}{corps}}. Démontrer un fait une seule fois, au niveau de la
structure, le démontre d'un coup pour tous les exemples.

\section{Lois de composition}

\begin{definition}\label{def:b1:structures:law}
Une \emph{loi de composition interne}\index{loi de composition interne}
sur un \omterm{def:b1:logic:sets}{ensemble} $E$ est une \omterm{def:b1:logic:map}{application} $E \times E \to E$, notée
$(x, y) \mapsto x * y$. Elle est \emph{associative} lorsque
$(x*y)*z = x*(y*z)$ toujours, \emph{commutative} lorsque $x * y = y * x$
toujours. Un élément $e$ est un \emph{élément neutre} lorsque
$e * x = x * e = x$ pour tout $x$~; alors $x'$ est un \emph{inverse} de
$x$ lorsque $x * x' = x' * x = e$.
\end{definition}

\begin{proposition}[Unicité]\label{prop:b1:structures:unique}
Une loi admet au plus un élément neutre~; pour une loi associative
possédant un élément neutre, chaque élément a au plus un inverse.
\end{proposition}

\begin{proof}
Si $e$ et $e'$ sont neutres~: $e = e * e' = e'$. Si $x'$ et $x''$
inversent $x$~: $x' = x' * e = x' * (x * x'') = (x' * x) * x'' = e * x''
= x''$.
\end{proof}

\section{Groupes}

\begin{definition}[Groupe]\label{def:b1:structures:group}
Un \emph{groupe}\index{groupe} $(G, *)$ est un \omterm{def:b1:logic:sets}{ensemble} muni d'une loi
associative admettant un élément neutre et dans lequel tout élément
possède un inverse. Le groupe est \emph{abélien}\index{groupe abélien}
lorsque la loi est commutative.
\end{definition}

\begin{example}\label{ex:b1:structures:groups}
$(\Z, +)$, $(\Q, +)$, $(\R, +)$, $(\C, +)$~; $(\Q^*, \times)$,
$(\R^*, \times)$, $(\C^*, \times)$, $(\mathbb{U}_n, \times)$ (\omterm{def:b1:complex:unity}{racines de
l'unité}, \cref{def:b1:complex:unity})~; l'\omterm{def:b1:logic:sets}{ensemble} $\mathfrak{S}(E)$ des
\omterm{def:b1:logic:inj}{bijections} d'un \omterm{def:b1:logic:sets}{ensemble} $E$ sur lui-même, muni de la composition --- le
\emph{groupe symétrique}\index{groupe symétrique} de $E$, non \omterm{def:b1:structures:group}{abélien} dès
que $\abs E \geq 3$. Ne sont pas des \omterm{def:b1:structures:group}{groupes}~: $(\N, +)$ (pas
d'inverses), $(\Z, \times)$ (seuls $\pm 1$ sont inversibles).
\end{example}

\begin{proposition}[Règles de calcul]\label{prop:b1:structures:rules}
Dans un \omterm{def:b1:structures:group}{groupe} $G$ (noté multiplicativement, de neutre $e$)~:
\begin{enumerate}
  \item simplification~: $ax = ay \implies x = y$ et $xa = ya \implies
        x = y$~;
  \item $(ab)^{-1} = b^{-1} a^{-1}$ et $(a^{-1})^{-1} = a$~;
  \item pour $a, b \in G$, chacune des équations $ax = b$ et $xa = b$
        admet une unique solution ($x = a^{-1}b$, resp.\ $x = b a^{-1}$).
\end{enumerate}
\end{proposition}

\begin{proof}
(1) On multiplie par $a^{-1}$ du bon côté, en utilisant
l'associativité. (2) $(b^{-1}a^{-1})(ab) = b^{-1}(a^{-1}a)b = b^{-1}b =
e$ et symétriquement~; l'unicité de l'inverse conclut~; le second
point est la \cref{prop:b1:structures:unique} appliquée à $a^{-1}$. (3) On
substitue, et (1) donne l'unicité.
\end{proof}

\begin{example}[Les symétries d'un rectangle]\label{ex:b1:structures:klein}
Un rectangle (non carré) admet exactement quatre isométries de lui-même
sur lui-même~: l'identité $e$, la symétrie d'axe horizontal $h$, la
symétrie d'axe vertical $v$, et le demi-tour $r$ autour du
centre. La composition fait de cet \omterm{def:b1:logic:sets}{ensemble} à quatre éléments un
\omterm{def:b1:structures:group}{groupe}~: chaque élément est son propre inverse ($h^2 = v^2 = r^2 = e$),
et le produit de deux éléments distincts autres que le neutre est le
troisième ($hv = vh = r$~: composer les deux symétries axiales donne le
demi-tour). La table est symétrique, donc le \omterm{def:b1:structures:group}{groupe} est \omterm{def:b1:structures:group}{abélien} --- et
pourtant ce n'est \emph{pas} le même \omterm{def:b1:structures:group}{groupe} que celui des rotations
$\mathbb U_4$ de l'\cref{ex:b1:structures:order}~: là, $\iu$ est d'\omterm{def:b1:structures:order}{ordre}
$4$, tandis qu'ici tout élément est d'\omterm{def:b1:structures:order}{ordre} $\leq 2$. Deux \omterm{def:b1:structures:group}{groupes} de
même taille peuvent donc avoir des structures multiplicatives
véritablement différentes --- la figure ci-dessous affiche les deux
tables côte à côte. Ce \omterm{def:b1:structures:group}{groupe} à quatre éléments réapparaîtra sous la
forme $\{\pm1\} \times \{\pm1\}$, et
l'\cref{exo:b1:structures:7} explique pourquoi tout \omterm{def:b1:structures:group}{groupe} dont tous les
carrés sont triviaux est, comme celui-ci, nécessairement \omterm{def:b1:structures:group}{abélien}.
\end{example}

\begin{omfigure}
\begin{tikzpicture}[scale=0.62]
  \foreach \i/\l in {0/$e$, 1/$\iu$, 2/$-1$, 3/$-\iu$} {
    \node[font=\small] at ({\i + 0.5}, 0.5) {\l};
    \node[font=\small] at (-0.5, {-\i - 0.5}) {\l};
  }
  \fill[omDef!25] (0,-1) rectangle ++(1,1);
  \fill[omDef!25] (3,-2) rectangle ++(1,1);
  \fill[omDef!25] (2,-3) rectangle ++(1,1);
  \fill[omDef!25] (1,-4) rectangle ++(1,1);
  \draw[gray] (0,-4) grid (4,0);
  \foreach \i/\j/\l in {0/0/$e$, 0/1/$\iu$, 0/2/$-1$, 0/3/$-\iu$,
    1/0/$\iu$, 1/1/$-1$, 1/2/$-\iu$, 1/3/$e$,
    2/0/$-1$, 2/1/$-\iu$, 2/2/$e$, 2/3/$\iu$,
    3/0/$-\iu$, 3/1/$e$, 3/2/$\iu$, 3/3/$-1$}
    \node[font=\scriptsize] at ({\j+0.5}, {-\i-0.5}) {\l};
  \begin{scope}[xshift=7.5cm]
    \foreach \i/\l in {0/$e$, 1/$h$, 2/$v$, 3/$r$} {
      \node[font=\small] at ({\i + 0.5}, 0.5) {\l};
      \node[font=\small] at (-0.5, {-\i - 0.5}) {\l};
    }
    \fill[omThm!25] (0,-1) rectangle ++(1,1);
    \fill[omThm!25] (1,-2) rectangle ++(1,1);
    \fill[omThm!25] (2,-3) rectangle ++(1,1);
    \fill[omThm!25] (3,-4) rectangle ++(1,1);
    \draw[gray] (0,-4) grid (4,0);
    \foreach \i/\j/\l in {0/0/$e$, 0/1/$h$, 0/2/$v$, 0/3/$r$,
      1/0/$h$, 1/1/$e$, 1/2/$r$, 1/3/$v$,
      2/0/$v$, 2/1/$r$, 2/2/$e$, 2/3/$h$,
      3/0/$r$, 3/1/$v$, 3/2/$h$, 3/3/$e$}
      \node[font=\scriptsize] at ({\j+0.5}, {-\i-0.5}) {\l};
  \end{scope}
\end{tikzpicture}

{\small Deux \omterm{def:b1:structures:group}{groupes} à quatre éléments~: $\mathbb U_4 = \{e, \iu,
-1, -\iu\}$ (à gauche) et le \omterm{def:b1:structures:group}{groupe} du rectangle (à droite), les
positions du neutre étant grisées. À gauche, le neutre serpente
(un élément d'\omterm{def:b1:structures:order}{ordre} $4$ engendre tout)~; à droite, il remplit la
diagonale (tout élément a pour carré $e$). Aucun changement de
noms ne transforme une table en l'autre~: les \omterm{def:b1:structures:group}{groupes} ne sont pas
isomorphes.}
\end{omfigure}

\begin{definition}[Sous-groupe]\label{def:b1:structures:subgroup}
Une partie $H$ d'un \omterm{def:b1:structures:group}{groupe} $G$ est un \emph{sous-groupe}\index{sous-groupe}
(on note $H \leq G$) lorsqu'elle contient $e$, est stable par la loi
et par passage à l'inverse. Alors $H$ est lui-même un \omterm{def:b1:structures:group}{groupe}.

\emph{Critère~:} une partie non vide $H \subseteq G$ est un \omterm{def:b1:structures:group}{sous-groupe}
si et seulement si
\[
\forall x, y \in H, \quad x y^{-1} \in H .
\]
\end{definition}

\begin{proof}[Démonstration du critère]
Un \omterm{def:b1:structures:subgroup}{sous-groupe} le vérifie de façon évidente. Réciproquement, soit
$H \neq \emptyset$ le vérifiant, et prenons $x_0 \in H$. Alors
$e = x_0 x_0^{-1} \in H$~; pour $y \in H$, $y^{-1} = e\,y^{-1} \in H$~;
et pour $x, y \in H$, $xy = x (y^{-1})^{-1} \in H$.
\end{proof}

\begin{example}\label{ex:b1:structures:subgroups}
$\mathbb{U}_n \leq (\C^*, \times)$~: non vide, et pour $z, w \in
\mathbb{U}_n$, $(zw^{-1})^n = z^n (w^n)^{-1} = 1$. Les \omterm{def:b1:structures:subgroup}{sous-groupes} de
$(\Z, +)$ sont exactement les $n\Z$ (démontré en
\cref{thm:b1:arith:gcd}). Une intersection de \omterm{def:b1:structures:subgroup}{sous-groupes} est toujours
un \omterm{def:b1:structures:subgroup}{sous-groupe}, mais une réunion ne l'est presque jamais
(\cref{exo:b1:structures:6}).
\end{example}

\begin{remark}[Pièges classiques avec les structures]\label{rem:b1:structures:pitfalls}
\begin{enumerate}
  \item \emph{La stabilité par la loi ne suffit pas.} $\N$ est
        stable par addition dans $\Z$ et contient $0$, mais n'est
        pas un \omterm{def:b1:structures:subgroup}{sous-groupe}~: les inverses manquent. Le critère
        $xy^{-1} \in H$ teste tout d'un coup --- mais seulement
        après avoir vérifié que $H \neq \emptyset$.
  \item \emph{Réflexes non \omterm{def:b1:structures:group}{abéliens}.} Dans un \omterm{def:b1:structures:group}{groupe} quelconque,
        $(ab)^2 = abab$, qui ne vaut $a^2b^2$ que si $a$ et $b$
        commutent~; de même $(ab)^{-1} = b^{-1}a^{-1}$, dans
        l'\omterm{def:b1:structures:order}{ordre} inverse. Toute identité importée de l'algèbre du
        lycée doit être redémontrée à partir des axiomes, ou
        signalée comme valable dans le cas commutatif.
  \item \emph{\omterm{def:b1:structures:morphism}{Noyau} contre image.} $\ker f$ vit dans l'\omterm{def:b1:logic:sets}{ensemble} de
        \emph{départ}, $\operatorname{im} f$ dans l'\omterm{def:b1:logic:sets}{ensemble}
        d'arrivée~; «~$f$ est \omterm{def:b1:logic:inj}{injectif} si et seulement si $\ker f$
        est trivial~» (\cref{prop:b1:structures:kernel}) n'a aucun
        analogue avec l'image ($\operatorname{im} f = G'$, c'est la
        \omterm{def:b1:logic:inj}{surjectivité}).
  \item \emph{Les anneaux ne sont pas des \omterm{def:b1:structures:group}{groupes} pour $\times$.}
        Dans un \omterm{def:b1:structures:ring}{anneau}, la plupart des éléments n'ont aucune raison
        d'être inversibles, et simplifier par $a$ exige que $a$
        soit inversible ou que l'\omterm{def:b1:structures:ring}{anneau} soit intègre~: dans
        $\Z/12\Z$, $\overline3\,\overline2 =
        \overline3\,\overline6$ alors que $\overline2 \neq
        \overline6$ (\cref{ex:b1:structures:zncomputation}).
\end{enumerate}
\end{remark}

\begin{definition}[Morphisme de groupes]\label{def:b1:structures:morphism}
Soient $(G, *)$ et $(G', \star)$ deux \omterm{def:b1:structures:group}{groupes}. Une \omterm{def:b1:logic:map}{application}
$f \colon G \to G'$ est un \emph{morphisme}\index{morphisme!de groupes}
lorsque
\[
\forall x, y \in G, \qquad f(x * y) = f(x) \star f(y).
\]
On a alors $f(e_G) = e_{G'}$ et $f(x^{-1}) = f(x)^{-1}$. Le
\emph{noyau}\index{noyau} et l'\emph{image} de $f$ sont
\[
\ker f = f^{-1}(\{e_{G'}\}) \leq G,
\qquad
\operatorname{im} f = f(G) \leq G' .
\]
Un morphisme \omterm{def:b1:logic:inj}{bijectif} est un \emph{isomorphisme}~; son \omterm{def:b1:logic:map}{application}
réciproque est alors automatiquement un morphisme.
\end{definition}

\begin{proof}[Démonstration des assertions]
$f(e) = f(e * e) = f(e)\star f(e)$, et en simplifiant par $f(e)$ on
obtient $e_{G'} = f(e)$. Ensuite $f(x)\star f(x^{-1}) = f(x x^{-1}) =
e_{G'}$ identifie $f(x^{-1})$ comme étant l'inverse. \omterm{def:b1:structures:morphism}{Noyau}~: $e \in
\ker f$~; si $x, y \in \ker f$, $f(xy^{-1}) = f(x)f(y)^{-1} = e$~; le
critère s'applique. Image~: même critère, avec $f(x)f(y)^{-1} =
f(xy^{-1})$. Réciproque d'un isomorphisme~: pour $u, v \in G'$,
écrivons $u = f(x)$, $v = f(y)$~; alors $f^{-1}(u \star v) =
f^{-1}(f(xy)) = xy = f^{-1}(u) f^{-1}(v)$.
\end{proof}

\begin{proposition}[Injectivité par le noyau]\label{prop:b1:structures:kernel}
Un \omterm{def:b1:structures:morphism}{morphisme} de \omterm{def:b1:structures:group}{groupes} $f$ est \omterm{def:b1:logic:inj}{injectif} si et seulement si $\ker f =
\{e\}$.
\end{proposition}

\begin{proof}
Si $f$ est \omterm{def:b1:logic:inj}{injectif}, $\ker f$ ne peut contenir que l'unique antécédent
de $e_{G'}$, à savoir $e$. Réciproquement, si $\ker f = \{e\}$ et
$f(x) = f(y)$, alors $f(xy^{-1}) = f(x) f(y)^{-1} = e_{G'}$, donc
$xy^{-1} = e$, c'est-à-dire $x = y$.
\end{proof}

\begin{example}\label{ex:b1:structures:expmorphism}
$\exp \colon (\R, +) \to (\R_+^*, \times)$ est un \omterm{def:b1:structures:morphism}{morphisme}
($\eu^{x+y} = \eu^x \eu^y$), \omterm{def:b1:logic:inj}{bijectif}
(\cref{prop:b1:functions:expln})~: les structures additive et
multiplicative sont isomorphes --- la raison d'être historique des
logarithmes. Autre \omterm{def:b1:structures:morphism}{morphisme}~: $\theta \mapsto \eu^{\iu\theta}$, de
$(\R, +)$ sur le cercle unité $(\mathbb{U}, \times)$, de \omterm{def:b1:structures:morphism}{noyau}
$2\pi\Z$.
\end{example}

\begin{example}[Le morphisme signe]\label{ex:b1:structures:signmorphism}
L'\omterm{def:b1:logic:map}{application} $s \colon (\R^*, \times) \to (\{\pm1\}, \times)$ qui
envoie $x$ sur son signe est un \omterm{def:b1:structures:morphism}{morphisme}~: le signe d'un produit
est le produit des signes. Son \omterm{def:b1:structures:morphism}{noyau} est $\intoo0{+\infty}$ (un
\omterm{def:b1:structures:subgroup}{sous-groupe}, comme le promet \cref{def:b1:structures:morphism}), son
image est $\{\pm1\}$ tout entier~: \omterm{def:b1:logic:inj}{surjectif}, massivement non
\omterm{def:b1:logic:inj}{injectif}. Deux leçons générales en miniature. D'abord, un \omterm{def:b1:structures:morphism}{morphisme}
peut écraser de l'information~: $s$ ne retient de $x$ qu'un seul
bit, et c'est là sa vertu --- les raisonnements de signe sont
exactement les calculs qui se factorisent par $s$. Ensuite, les
\omterm{def:b1:structures:morphism}{morphismes} à valeurs dans $\{\pm1\}$ sont les «~invariants~» les
plus simples~: la signature des \omterm{def:b1:counting:objects}{permutations}, construite dans le
devoir maison de ce chapitre, est le même phénomène sur le \omterm{def:b1:structures:group}{groupe}
$\mathfrak S_n$, et tous les raisonnements de parité qu'elle
alimente passent par un tel \omterm{def:b1:structures:morphism}{morphisme} à deux valeurs.
\end{example}

\begin{definition}[Puissances, ordre d'un élément]\label{def:b1:structures:order}
Dans un \omterm{def:b1:structures:group}{groupe} $G$ (en notation multiplicative), on pose $x^0 = e$,
$x^{k+1} = x^k x$ et $x^{-k} = (x^k)^{-1}$ pour $k \in \N$~; on a alors
$x^{k+l} = x^k x^l$ pour tous $k, l \in \Z$, de sorte que $k \mapsto
x^k$ est un \omterm{def:b1:structures:morphism}{morphisme} $(\Z, +) \to G$ dont l'image $\langle x \rangle =
\{x^k : k \in \Z\}$ est un \omterm{def:b1:structures:subgroup}{sous-groupe}, le \omterm{def:b1:structures:subgroup}{sous-groupe}
\emph{engendré} par $x$. L'\emph{ordre}\index{ordre d'un élément} de
$x$ est le plus petit $m \geq 1$ tel que $x^m = e$, s'il en existe un
(alors $\langle x\rangle = \{e, x, \dots, x^{m-1}\}$ a exactement $m$
éléments, et $x^k = e \iff m \mid k$), et $\infty$ sinon.
\end{definition}

\begin{example}\label{ex:b1:structures:order}
Dans $(\C^*, \times)$~: $\iu$ est d'\omterm{def:b1:structures:order}{ordre} $4$, avec $\langle \iu \rangle
= \{1, \iu, -1, -\iu\} = \mathbb{U}_4$~; plus généralement $\omega =
\eu^{2\iu\pi/n}$ est d'\omterm{def:b1:structures:order}{ordre} $n$ et $\langle\omega\rangle =
\mathbb{U}_n$. Dans $(\Z, +)$, tout $x \neq 0$ est d'\omterm{def:b1:structures:order}{ordre} infini.
Pourquoi les affirmations de la définition sont vraies~: si $x$ est
d'\omterm{def:b1:structures:order}{ordre} $m$, effectuons la division euclidienne d'un $k$ quelconque par
$m$ ($k = mq + r$, $0 \leq r < m$,
\cref{thm:b1:arith:division})~: $x^k = (x^m)^q x^r = x^r$, donc les
puissances sont périodiques de période $m$, les éléments énumérés sont
deux à deux distincts par minimalité de $m$, et $x^k = e$ force $r = 0$.
Les \omterm{def:b1:structures:order}{ordres} des \omterm{def:b1:counting:objects}{permutations} sont calculés dans le devoir maison
ci-dessous.
\end{example}

\begin{example}[Ordres dans $\mathbb U_{12}$]\label{ex:b1:structures:orderU12}
Quel est l'\omterm{def:b1:structures:order}{ordre} de $\omega^k$ dans $\mathbb U_n$, pour $\omega =
\eu^{2\iu\pi/n}$~? On a $(\omega^k)^m = 1$ si et seulement si
$n \mid km$, et en écrivant $d = \gcd(n, k)$, $n = dn'$, $k = dk'$
avec $\gcd(n', k') = 1$~: $n \mid km \iff n' \mid k'm \iff n' \mid m$
(\omterm{thm:b1:arith:gauss}{lemme de Gauss}, \cref{thm:b1:arith:gauss}). Le plus petit tel
$m \geq 1$ est $n' = \frac{n}{\gcd(n,k)}$. Dans $\mathbb U_{12}$ par
exemple, $\omega^8$ est d'\omterm{def:b1:structures:order}{ordre} $\frac{12}{\gcd(12,8)} = 3$ (en effet
$\omega^8 = \eu^{4\iu\pi/3} \in \mathbb U_3$), tandis que $\omega^5$
est d'\omterm{def:b1:structures:order}{ordre} $12$~: il engendre le \omterm{def:b1:structures:group}{groupe} tout entier, bien qu'il ne
soit pas le générateur «~standard~». Compter les générateurs --- les
$k$ tels que $\gcd(k, n) = 1$ --- redonne les dénombrements d'entiers
premiers avec $n$ de l'\cref{ex:b1:counting:coprime}~: théorie des
\omterm{def:b1:structures:group}{groupes} et dénombrement se rejoignent.
\end{example}

\section{Anneaux et corps}

\begin{definition}[Anneau]\label{def:b1:structures:ring}
Un \emph{anneau}\index{anneau} $(A, +, \times)$ est un \omterm{def:b1:logic:sets}{ensemble} muni de
deux lois telles que~: $(A, +)$ est un \omterm{def:b1:structures:group}{groupe abélien} (de neutre $0$)~;
$\times$ est associative et possède un élément neutre $1$~; et $\times$
est distributive par rapport à $+$ des deux côtés. L'anneau est
\emph{commutatif} lorsque $\times$ l'est. Un élément $a$ est
\emph{inversible} (on dit aussi que c'est une \emph{unité}) lorsque
$ab = ba = 1$ pour un certain $b$~; les inversibles forment un \omterm{def:b1:structures:group}{groupe}
$(A^\times, \times)$.
\end{definition}

\begin{proof}[Démonstration~: les inversibles forment un groupe]
Stabilité~: si $a, a'$ sont inversibles, d'inverses $b, b'$, alors
\[
(aa')(b'b) = a(a'b')b = a\,1\,b = ab = 1,
\qquad (b'b)(aa') = 1
\]
symétriquement, donc $aa'$ est inversible. L'élément $1$ est
inversible (il est son propre inverse), l'associativité est héritée
de $A$, et l'inverse $b$ d'un inversible $a$ est lui-même inversible
(d'inverse $a$). Ainsi $(A^\times, \times)$ vérifie tous les axiomes
de \omterm{def:b1:structures:group}{groupe}. Tous les \omterm{def:b1:structures:group}{groupes} de ce livre qui ne sont pas construits à
partir de \omterm{def:b1:counting:objects}{permutations} apparaissent de cette façon~:
$\Q^* = \Q^\times$, $\R^*$, $\C^*$, les inversibles de $\Z/n\Z$
ci-dessous, et plus tard les matrices inversibles
(\cref{ch:b1:matrices}).
\end{proof}

\begin{example}\label{ex:b1:structures:rings}
$\Z, \Q, \R, \C$ sont des anneaux commutatifs~; $\Z^\times = \{1, -1\}$,
$\Q^\times = \Q^*$. Plus tard~: les anneaux de polynômes $K[X]$
(\cref{ch:b1:poly}), les anneaux de matrices (non commutatifs,
\cref{ch:b1:matrices}), et $\Z/n\Z$ ci-dessous. Dans tout \omterm{def:b1:structures:ring}{anneau},
$0 \times a = 0$ (par distributivité~: $0a = (0+0)a = 0a + 0a$), et
$(-1)a = -a$.
\end{example}

\begin{example}[Idempotents~: des phénomènes nouveaux dans des anneaux nouveaux]\label{ex:b1:structures:idempotents}
Dans $\Z$, l'équation $x^2 = x$, c'est-à-dire $x(x - 1) = 0$, n'a que
les solutions $0$ et $1$. Dans $\Z/6\Z$, en testant toutes les
classes~: $\overline0^2 = \overline0$, $\overline1^2 = \overline1$,
$\overline3^2 = \overline9 = \overline3$ et $\overline4^2 =
\overline{16} = \overline4$ --- \emph{quatre} idempotents. Les deux
exotiques proviennent de diviseurs de zéro~: $\overline3\,(\overline3 -
\overline1) = \overline3 \times \overline2 = \overline6 =
\overline0$, sans qu'aucun des deux facteurs ne soit nul. De tels
calculs recalibrent l'intuition~: les faits familiers sur les
équations survivent dans les anneaux intègres et dans les \omterm{def:b1:structures:field}{corps},
mais un \omterm{def:b1:structures:ring}{anneau} quelconque peut se comporter --- et se comporte ---
tout autrement~; voir aussi les anneaux de Boole de
l'\cref{exo:b1:structures:10}, où \emph{tout} élément est idempotent.
\end{example}

\begin{proposition}[Formule du binôme dans un anneau commutatif]\label{prop:b1:structures:binomial}
Si $a, b$ sont des éléments d'un \omterm{def:b1:structures:ring}{anneau} commutatif (plus généralement,
si $ab = ba$), alors pour $n \in \N$~:
\[
(a+b)^n = \sum_{k=0}^n \binom nk a^k b^{n-k},
\qquad
a^n - b^n = (a - b) \sum_{k=0}^{n-1} a^k b^{\,n-1-k} .
\]
\end{proposition}

\begin{proof}
Les démonstrations du \cref{thm:b1:counting:binomial} et de l'identité
géométrique n'utilisent que l'associativité, la commutativité des deux
éléments et la distributivité --- elles s'appliquent mot pour mot.
\end{proof}

\begin{example}[La formule du binôme dans un anneau inhabituel]\label{ex:b1:structures:binomialring}
Deux bénéfices immédiats de cette généralité. Dans $\Z/p\Z$ ($p$
premier), les coefficients binomiaux intermédiaires sont nuls (c'est
la première étape du \cref{thm:b1:arith:fermat}), de sorte que la
formule s'effondre en le \emph{rêve du débutant}
\[
(a + b)^p = a^p + b^p \qquad \text{dans } \Z/p\Z ,
\]
une véritable identité dans cet \omterm{def:b1:structures:ring}{anneau}, si criminelle qu'elle paraisse
sur $\R$.
Et dans tout \omterm{def:b1:structures:ring}{anneau} commutatif contenant un élément $\varepsilon$
tel que $\varepsilon^2 = 0$, la formule se tronque~:
$(a + \varepsilon)^n = a^n + n\,a^{n-1}\varepsilon$, tous les termes
suivants portant un facteur $\varepsilon^2 = 0$. Le coefficient
$n\,a^{n-1}$ de $\varepsilon$ est la dérivée de $x^n$ --- ce n'est
pas un hasard, et c'est un premier indice que les dérivées relèvent
autant de l'algèbre que de l'analyse (comparer avec la dérivée
formelle du \cref{ch:b1:poly}).
\end{example}

\begin{definition}[Anneau intègre, corps]\label{def:b1:structures:field}
Un \omterm{def:b1:structures:ring}{anneau} commutatif $A \neq \{0\}$ est un \emph{anneau
intègre}\index{anneau intègre} lorsqu'il n'a pas de diviseur de zéro~:
$ab = 0 \implies a = 0$ ou $b = 0$. C'est un \emph{corps}\index{corps}
lorsque tout élément non nul est inversible. Tout corps est intègre
($ab = 0$ et $a \neq 0$ donnent $b = a^{-1}ab = 0$).
\end{definition}

\begin{example}
$\Q$, $\R$, $\C$ sont des \omterm{def:b1:structures:field}{corps}~; $\Z$ est un \omterm{def:b1:structures:field}{anneau intègre} mais n'est
pas un \omterm{def:b1:structures:field}{corps}. Dans un \omterm{def:b1:structures:field}{anneau intègre}, on peut simplifier pour $\times$~:
$ab = ac$ et $a \neq 0$ entraînent $b = c$.
\end{example}

\section{L'anneau $\Z/n\Z$}

\begin{definition}\label{def:b1:structures:zn}
Fixons $n \in \N^*$. Les classes de \omterm{def:b1:arith:congruence}{congruence} modulo $n$
(\cref{ex:b1:logic:congruence}) forment un \omterm{def:b1:logic:sets}{ensemble}
$\Z/n\Z$\index{Z/nZ@$\Z/n\Z$} à $n$ éléments, notées $\overline 0,
\overline 1, \dots, \overline{n-1}$. Les opérations
\[
\overline a + \overline b = \overline{a + b},
\qquad
\overline a \times \overline b = \overline{ab}
\]
sont bien définies --- les classes des résultats ne dépendent pas des
représentants, précisément parce que la \omterm{def:b1:arith:congruence}{congruence} est compatible avec
$+$ et $\times$ (\cref{def:b1:arith:congruence}) --- et font de
$\Z/n\Z$ un \omterm{def:b1:structures:ring}{anneau} commutatif.
\end{definition}

\begin{theorem}[Inversibles de $\Z/n\Z$~; les corps $\Z/p\Z$]\label{thm:b1:structures:zp}
\begin{enumerate}
  \item $\overline a$ est inversible dans $\Z/n\Z$ si et seulement si
        $\gcd(a, n) = 1$.
  \item $\Z/n\Z$ est un \omterm{def:b1:structures:field}{corps} si et seulement si $n$ est premier.
\end{enumerate}
\end{theorem}

\begin{proof}
(1) n'est autre que la \cref{prop:b1:arith:invmod} réécrite en termes de
classes.

(2) Si $n = p$ est premier, tout $\overline a \neq \overline 0$ vérifie
$p \nmid a$, donc $\gcd(a, p) = 1$~: il est inversible par (1) --- c'est
un \omterm{def:b1:structures:field}{corps}. Si $n = ab$ avec $1 < a, b < n$, alors $\overline a\,
\overline b = \overline n = \overline 0$ avec $\overline a, \overline b
\neq \overline 0$~: il y a des diviseurs de zéro, donc l'\omterm{def:b1:structures:ring}{anneau} n'est
même pas intègre~; et $n = 1$ donne l'\omterm{def:b1:structures:ring}{anneau} nul, exclu.
\end{proof}

\begin{example}[Combien de racines carrées de $1$~?]\label{ex:b1:structures:sqrtone}
Résolvons $x^2 = \overline 1$ dans $\Z/8\Z$ puis dans $\Z/7\Z$. En
testant les huit classes modulo $8$~: $1^2 = 1$, $3^2 = 9 \equiv 1$,
$5^2 = 25 \equiv 1$, $7^2 = 49 \equiv 1$ --- \emph{quatre} solutions
$\{\overline1, \overline3, \overline5, \overline7\}$, alors même que
le polynôme $X^2 - 1$ est de degré $2$. Dans le \omterm{def:b1:structures:field}{corps} $\Z/7\Z$,
en revanche, $x^2 = \overline1$ signifie $(x - \overline1)(x +
\overline1) = \overline0$, et un \omterm{def:b1:structures:field}{corps} n'a pas de diviseur de zéro~:
$x = \pm\overline1$, deux solutions seulement. L'échec modulo $8$
s'explique~: $(3-1)(3+1) = 2 \times 4 = 8 \equiv 0$ sans qu'aucun des
deux facteurs ne s'annule. Morale~: la règle familière «~une équation
de degré $d$ a au plus $d$ racines~» est un théorème sur les
\emph{anneaux intègres} (le \cref{cor:b1:poly:nroots} le démontre sur un
\omterm{def:b1:structures:field}{corps})~; dans un \omterm{def:b1:structures:ring}{anneau} à diviseurs de zéro, elle tombe en défaut sans
prévenir ---
et c'est exactement pourquoi la démonstration par appariement du
théorème de Wilson (\cref{exo:b1:arith:11}) exigeait $p$ premier.
\end{example}

\begin{example}[Calculer dans $\Z/n\Z$]\label{ex:b1:structures:zncomputation}
Dans $\Z/12\Z$~: les inversibles sont $\overline 1, \overline 5,
\overline 7, \overline{11}$ (les classes premières avec $12$), et chacun
est son propre inverse ($5^2 = 25 \equiv 1$, $7^2 = 49 \equiv 1$,
$11^2 = 121 \equiv 1$). L'équation $\overline 3\, x = \overline 6$ a
\emph{trois} solutions ($x \in \{\overline 2, \overline 6,
\overline{10}\}$)~: sans inversibilité, pas de simplification. Dans
$\Z/11\Z$ en revanche, toute équation $\overline a x = \overline b$ avec
$\overline a \neq \overline 0$ a exactement une solution.
\end{example}

\begin{example}[Les axiomes de groupe comme permis de résoudre]\label{ex:b1:structures:solving}
Dans le \omterm{def:b1:structures:group}{groupe} $\bigl((\Z/7\Z)^*, \times\bigr)$, résolvons $\overline
3\,x = \overline 5$. D'après la \cref{prop:b1:structures:rules}~(3), la
solution existe, elle est unique, et vaut $\overline3^{-1}\,
\overline5$~; comme $\overline3 \times \overline5 = \overline{15}
= \overline1$, l'inverse de $\overline 3$ est $\overline 5$, d'où
\[
x = \overline5 \times \overline5 = \overline{25} = \overline4,
\qquad\text{vérification~: } \overline3 \times \overline4 =
\overline{12} = \overline5 .
\]
Ce qui compte, moins que la réponse, c'est la garantie~: dans un
\omterm{def:b1:structures:group}{groupe}, toute équation de ce type admet une solution unique
\emph{avant} tout calcul, si bien qu'une procédure de résolution ne
peut jamais tomber sur «~pas de solution~» ou «~plusieurs~».
Comparer avec $\overline3\,x = \overline6$ dans $\Z/12\Z$ ci-dessus,
où la garantie tombe en défaut --- savoir dans quelle structure on
travaille, c'est savoir ce que l'on peut tenir pour acquis.
\end{example}

\begin{example}[Produits directs]\label{ex:b1:structures:product}
Si $G$ et $H$ sont des \omterm{def:b1:structures:group}{groupes}, l'\omterm{def:b1:logic:sets}{ensemble} produit $G \times H$ muni
de la loi composante par composante $(g, h)(g', h') = (gg', hh')$ est
un \omterm{def:b1:structures:group}{groupe}~: les axiomes se vérifient coordonnée par coordonnée, avec
pour neutre $(e_G, e_H)$ et pour inverses $(g^{-1}, h^{-1})$. Les
\omterm{def:b1:structures:order}{ordres} se combinent par le PPCM~: $(g, h)^m = (g^m, h^m)$ est le
neutre si et seulement si l'\omterm{def:b1:structures:order}{ordre} de $g$ et l'\omterm{def:b1:structures:order}{ordre} de $h$ \omterm{def:b1:arith:divides}{divisent}
tous deux $m$. Ainsi, dans $\Z/2\Z \times \Z/2\Z$ (noté additivement),
tout élément non nul est d'\omterm{def:b1:structures:order}{ordre} $2$ --- c'est exactement le \omterm{def:b1:structures:group}{groupe} du
rectangle de l'\cref{ex:b1:structures:klein} en coordonnées --- alors
que $\Z/4\Z$ possède un élément d'\omterm{def:b1:structures:order}{ordre} $4$~: seconde démonstration,
sans le moindre calcul, du fait que les deux \omterm{def:b1:structures:group}{groupes} de \omterm{def:b1:counting:card}{cardinal} $4$
ne sont pas isomorphes (un isomorphisme conserve les \omterm{def:b1:structures:order}{ordres}). Les
produits sont le moyen le plus simple de fabriquer des \omterm{def:b1:structures:group}{groupes}
nouveaux à partir d'anciens, et le plan $\R^2 = \R \times \R$ du
\cref{ch:b1:vspaces} est l'exemple le plus important de cette
construction.
\end{example}

\begin{remark}[Fermat, structurellement]
Dans le \omterm{def:b1:structures:field}{corps} $\Z/p\Z$, les classes non nulles forment un \omterm{def:b1:structures:group}{groupe}
multiplicatif à $p - 1$ éléments, et le \omterm{thm:b1:arith:fermat}{petit théorème de Fermat}
(\cref{thm:b1:arith:fermat}) affirme~: tout élément $x$ de ce \omterm{def:b1:structures:group}{groupe}
vérifie $x^{p-1} = \overline 1$. C'est un cas particulier d'un fait
général sur les \omterm{def:b1:structures:group}{groupes} finis (le théorème de Lagrange), démontré en
deuxième année~; la démonstration par appariement du théorème de
Wilson (\cref{exo:b1:arith:11}) avait déjà cette saveur de théorie des
\omterm{def:b1:structures:group}{groupes}.
\end{remark}

\begin{remark}[Interlude~: ce que rapporte l'abstraction]\label{rem:b1:structures:interlude}
On peut légitimement demander ce que l'on a gagné à démontrer,
disons, la \cref{prop:b1:structures:unique} pour une loi abstraite
plutôt que pour des nombres. La réponse est un effet de levier. Cet
argument de deux lignes couvre désormais, d'un seul coup~: les
réciproques des fonctions pour la composition
(\cref{thm:b1:logic:inverse}, dont il répète mot pour mot la
démonstration d'unicité), les inverses modulo $n$
(\cref{prop:b1:arith:invmod}), les inverses des réels non nuls, ceux
des inversibles d'un \omterm{def:b1:structures:ring}{anneau} quelconque et --- sans même les avoir
vues --- ceux des matrices inversibles du \cref{ch:b1:matrices}, où
l'unicité de $A^{-1}$ ne demandera pas une seule ligne de
démonstration. La même économie vaut pour
la \cref{prop:b1:structures:kernel} (un unique critère d'\omterm{def:b1:logic:inj}{injectivité},
réutilisé pour les \omterm{def:b1:logic:map}{applications} linéaires au \cref{ch:b1:linmaps}) et
pour le critère de \omterm{def:b1:structures:subgroup}{sous-groupe}. L'abstraction n'est pas ici la
généralité pour elle-même~: c'est le refus de démontrer cinq fois le
même lemme sous cinq noms différents. Le prix --- garder trace des
axiomes que chaque énoncé a réellement utilisés --- est précisément
ce que les exercices de ce chapitre entraînent.
\end{remark}

\begin{remark}[Où ce chapitre sert]\label{rem:b1:structures:whereused}
Le vocabulaire de ce chapitre est la grammaire de tout le reste du
volume. Anneaux et \omterm{def:b1:structures:field}{corps} organisent le \cref{ch:b1:poly} ($K[X]$ est
un \omterm{def:b1:structures:ring}{anneau} qui imite $\Z$) et le \cref{ch:b1:fractions} ($K(X)$ en est
le \omterm{def:b1:structures:field}{corps} des fractions)~; les espaces vectoriels
(\cref{ch:b1:vspaces}) sont des \omterm{def:b1:structures:group}{groupes abéliens} sur lesquels agit un
\omterm{def:b1:structures:field}{corps}~; les matrices (\cref{ch:b1:matrices}) forment le premier \omterm{def:b1:structures:ring}{anneau}
sérieusement non commutatif du volume, et leurs éléments inversibles
un \omterm{def:b1:structures:group}{groupe} dont l'étude est l'algèbre linéaire elle-même. \omterm{def:b1:structures:morphism}{Morphismes} et
noyaux reviennent sous la forme des \omterm{def:b1:logic:map}{applications} linéaires et de leurs
noyaux au \cref{ch:b1:linmaps} --- la \cref{prop:b1:structures:kernel}
\emph{est} le critère d'\omterm{def:b1:logic:inj}{injectivité} de ce chapitre-là, démontré ici
une fois pour toutes. Le \omterm{ex:b1:structures:groups}{groupe symétrique}, vedette du devoir maison
ci-dessous, fournit la signature sur laquelle sont bâtis les
déterminants au \cref{ch:b1:det}.
\end{remark}

\section{Exercices}

\begin{exercise}[$\star$]\label{exo:b1:structures:1}
Sur $E = \R \setminus \{1\}$, on définit $x * y = x + y - xy$. Démontrer
que $(E, *)$ est un \omterm{def:b1:structures:group}{groupe abélien}. \emph{(Identifier l'élément neutre
et l'inverse de $x$~; vérifier la stabilité~: pourquoi a-t-on
$x * y \neq 1$~?)}
\end{exercise}

\begin{exercise}[$\star$]\label{exo:b1:structures:2}
Lesquels des \omterm{def:b1:logic:sets}{ensembles} suivants sont des \omterm{def:b1:structures:group}{groupes}~?
\begin{enumerate}
  \item $(\intoo{0}{+\infty}, \times)$~;
  \item $(\{-1, 0, 1\}, +)$~;
  \item $(\Q^*, \times)$~;
  \item l'\omterm{def:b1:logic:sets}{ensemble} des entiers impairs muni de l'addition.
\end{enumerate}
\end{exercise}

\begin{exercise}[$\star$]\label{exo:b1:structures:3}
Écrire la table de composition du \omterm{ex:b1:structures:groups}{groupe symétrique} $\mathfrak{S}_3$
de $\{1,2,3\}$ (six \omterm{def:b1:logic:inj}{bijections}~: l'identité, trois transpositions, deux
$3$-cycles), et exhiber deux éléments qui ne commutent pas.
\end{exercise}

\begin{exercise}[$\star$]\label{exo:b1:structures:4}
Démontrer que $H = \{z \in \C^* : \abs z = 1\}$ est un \omterm{def:b1:structures:subgroup}{sous-groupe} de
$(\C^*, \times)$, et que $\R_+^*$ en est un autre~; $H \cup \R_+^*$
est-il un \omterm{def:b1:structures:subgroup}{sous-groupe}~?
\end{exercise}

\begin{exercise}[$\star\star$]\label{exo:b1:structures:5}
Soit $f \colon (\R, +) \to (\C^*, \times)$, $\theta \mapsto
\eu^{\iu\theta}$. Démontrer que $f$ est un \omterm{def:b1:structures:morphism}{morphisme}, calculer $\ker f$
et $\operatorname{im} f$, et déduire de
la \cref{prop:b1:structures:kernel} que $f$ n'est pas \omterm{def:b1:logic:inj}{injectif}. Restreindre
l'\omterm{def:b1:logic:sets}{ensemble} de départ de façon à rendre $f$ \omterm{def:b1:logic:inj}{injectif} sur un intervalle
aussi grand que possible.
\end{exercise}

\begin{exercise}[$\star\star$]\label{exo:b1:structures:6}
Soient $H, K$ des \omterm{def:b1:structures:subgroup}{sous-groupes} de $G$. Démontrer que $H \cap K$ est un
\omterm{def:b1:structures:subgroup}{sous-groupe}, et que $H \cup K$ n'est un \omterm{def:b1:structures:subgroup}{sous-groupe} \emph{que} si
$H \subseteq K$ ou $K \subseteq H$. \emph{(Si $h \in H \setminus K$ et
$k \in K \setminus H$, où $hk$ peut-il bien vivre~?)}
\end{exercise}

\begin{exercise}[$\star\star$]\label{exo:b1:structures:7}
Un \omterm{def:b1:structures:group}{groupe} $G$ vérifie $x^2 = e$ pour tout $x \in G$. Démontrer que $G$
est \omterm{def:b1:structures:group}{abélien}. \emph{(Développer $(xy)^2$.)}
\end{exercise}

\begin{exercise}[$\star\star$]\label{exo:b1:structures:8}
Dans $\Z/18\Z$~: énumérer les inversibles et trouver l'inverse de
$\overline 5$~; résoudre $\overline 5\, x = \overline 7$~; résoudre
$\overline 6\, x = \overline 3$ et $\overline 6\, x = \overline{12}$.
\end{exercise}

\begin{exercise}[$\star\star$]\label{exo:b1:structures:9}
Démontrer que l'\omterm{def:b1:logic:sets}{ensemble} $\Z[\sqrt 2] = \{a + b\sqrt 2 : a, b \in \Z\}$
est un \omterm{def:b1:structures:ring}{anneau} (un \omterm{def:b1:structures:ring}{sous-anneau} de $\R$), et que $1 + \sqrt 2$ en est un
inversible dont les puissances sont deux à deux distinctes --- de sorte
que $\Z[\sqrt 2]^\times$ est infini, contrairement à $\Z^\times$.
\end{exercise}

\begin{exercise}[$\star\star\star$]\label{exo:b1:structures:10}
(Anneaux de Boole) Soit $A$ un \omterm{def:b1:structures:ring}{anneau} dans lequel $x^2 = x$ pour tout
$x$. Démontrer que $x + x = 0$ pour tout $x$, et que $A$ est commutatif.
\emph{(Développer $(x+x)^2$ et $(x+y)^2$.)} Donner un exemple d'un tel
\omterm{def:b1:structures:ring}{anneau} à partir de $\mathcal{P}(E)$, en prenant la différence symétrique
pour addition et l'intersection pour multiplication.
\end{exercise}

\begin{exercise}[$\star\star\star$]\label{exo:b1:structures:11}
Soit $G$ un \omterm{def:b1:structures:group}{groupe} dans lequel, pour un certain $n \geq 1$ fixé,
$(xy)^n = x^n y^n$, $(xy)^{n+1} = x^{n+1}y^{n+1}$ et $(xy)^{n+2} =
x^{n+2}y^{n+2}$ pour tous $x, y$. Démontrer que $G$ est \omterm{def:b1:structures:group}{abélien}.
\emph{(À partir des trois identités, établir d'abord $y^n x = x y^n$,
puis $y^{n+1} x = x y^{n+1}$, et conclure.)}
\end{exercise}

\begin{exercise}[$\star\star$]\label{exo:b1:structures:12}
\begin{enumerate}
  \item Déterminer tous les \omterm{def:b1:structures:morphism}{morphismes} de \omterm{def:b1:structures:group}{groupes} de $(\Z, +)$ dans
        $(\Z, +)$.
  \item Démontrer que le seul \omterm{def:b1:structures:morphism}{morphisme} de \omterm{def:b1:structures:group}{groupes} de $(\Q, +)$ dans
        $(\Z, +)$ est le \omterm{def:b1:structures:morphism}{morphisme} nul. \emph{(Pour $x \in \Q$ et
        $n \in \N^*$, comparer $f(x)$ et $n\,f(x/n)$.)}
\end{enumerate}
\end{exercise}

\section{Problème~: le groupe symétrique et le taquin}

\begin{problem}\label{pb:b1:structures:1}
Le \omterm{def:b1:structures:group}{groupe} $\mathfrak S_n$ des \omterm{def:b1:counting:objects}{permutations} de $\intint1n$ est le plus
ancien \omterm{def:b1:structures:group}{groupe} des mathématiques, et il reste le plus instructif. Ce
problème en construit la théorie de A à Z --- cycles, engendrement
par les transpositions, le \omterm{def:b1:structures:morphism}{morphisme} \emph{signature}
$\varepsilon \colon \mathfrak S_n \to \{\pm1\}$ (dont l'existence n'a
rien d'évident), et le \omterm{pb:b1:structures:1}{groupe alterné} $\mathfrak A_n$ engendré par les
$3$-cycles --- puis la met à profit sur un casse-tête classique~:
dans le jeu de taquin $3 \times 3$, aucune suite de coups ne peut
échanger deux pièces en laissant tout le reste en place.
\index{groupe symétrique}\index{signature}\index{groupe alterné}
Les \omterm{def:b1:counting:objects}{permutations} agissent sur $\intint1n$~; le produit $\sigma\tau$
signifie «~appliquer $\tau$ d'abord~»~; $[\,v_1, \dots, v_n]$ désigne
la \omterm{def:b1:counting:objects}{permutation} envoyant $i$ sur $v_i$.

\medskip
\noindent\textbf{Partie I --- Cycles et transpositions.}
\begin{enumerate}
  \item Justifier $\abs{\mathfrak S_n} = n!$
        (\cref{thm:b1:counting:counts}). Dans $\mathfrak S_3$,
        calculer les deux produits de $\sigma = [2, 3, 1]$ et $\tau =
        [1, 3, 2]$, et en conclure que $\mathfrak S_3$ n'est pas
        \omterm{def:b1:structures:group}{abélien}.
  \item Un \emph{$k$-cycle} $(a_1\ a_2\ \dots\ a_k)$ ($k \geq 2$,
        les $a_i$ deux à deux distincts) envoie $a_1 \mapsto a_2
        \mapsto \dots \mapsto a_k \mapsto a_1$ et fixe tout le
        reste~; son \emph{support} est $\{a_1, \dots, a_k\}$.
        Démontrer que deux cycles de supports disjoints commutent.
  \item Démontrer que tout $\sigma \in \mathfrak S_n$ est un produit
        de cycles à supports deux à deux disjoints, et que cette
        décomposition est unique à l'\omterm{def:b1:structures:order}{ordre} des facteurs près.
        \emph{(Considérer, pour chaque $i$, la suite $i, \sigma(i),
        \sigma^2(i), \dots$~: elle revient nécessairement à $i$~; les
        \emph{orbites} obtenues forment une \omterm{thm:b1:logic:partition}{partition} de
        $\intint1n$, et $\sigma$ agit sur chacune comme un cycle.)}
  \item Décomposer $\sigma = [4, 1, 5, 2, 3, 7, 8, 6] \in
        \mathfrak S_8$ en cycles disjoints. En définissant
        l'\emph{\omterm{def:b1:structures:order}{ordre}} de $\sigma$ comme dans
        la \cref{def:b1:structures:order}, démontrer que l'\omterm{def:b1:structures:order}{ordre} d'un
        produit de cycles disjoints est le PPCM de leurs longueurs,
        et calculer l'\omterm{def:b1:structures:order}{ordre} de ce $\sigma$.
  \item Démontrer l'identité télescopique
        \[
        (a_1\ a_2\ \dots\ a_k)
        = (a_1\ a_k)(a_1\ a_{k-1})\cdots(a_1\ a_2) ,
        \]
        et en conclure que toute \omterm{def:b1:counting:objects}{permutation} est un produit de
        transpositions. Écrire le $\sigma$ de la question 4 sous cette
        forme.
  \item Montrer de plus que les transpositions \emph{adjacentes}
        $(i\ \ i{+}1)$ suffisent~: pour $a < b$,
        \[
        (a\ b) = (a\ \ a{+}1)(a{+}1\ \ a{+}2)\cdots(b{-}1\ \ b)
        \cdots(a{+}1\ \ a{+}2)(a\ \ a{+}1),
        \]
        soit un produit de $2(b - a) - 1$ transpositions adjacentes ---
        un nombre \emph{impair} (cette parité servira deux fois plus
        bas).
\end{enumerate}

\medskip
\noindent\textbf{Partie II --- La signature existe.} Pour $\sigma
\in \mathfrak S_n$, soit
\[
N(\sigma) = \#\bigl\{(i, j) : i < j,\ \sigma(i) >
\sigma(j)\bigr\}
\]
son nombre d'\emph{inversions}, et posons $\varepsilon(\sigma) =
(-1)^{N(\sigma)}$.
\begin{enumerate}[resume]
  \item Calculer $N$ et $\varepsilon$ pour l'identité, pour une
        transposition $(i\ \ i{+}1)$, et pour $[2, 3, 1]$.
  \item Démontrer que pour tout $\sigma$ et toute transposition
        adjacente $\tau = (i\ \ i{+}1)$~:
        $N(\sigma\tau) = N(\sigma) \pm 1$. \emph{(Composer à droite
        par $\tau$ échange les valeurs situées aux positions $i$ et
        $i + 1$~; exactement une paire change de statut
        d'inversion.)}
  \item En déduire, à l'aide de la question 6, que pour
        \emph{toute} transposition $\tau$, $\varepsilon(\sigma\tau) =
        -\varepsilon(\sigma)$~; conclure que si $\sigma$ est un
        produit de $p$ transpositions, alors $\varepsilon(\sigma)
        = (-1)^p$ --- en particulier la parité de $p$ ne dépend que
        de $\sigma$, et non de la factorisation choisie --- et que
        $\varepsilon \colon \mathfrak S_n \to \{\pm 1\}$ est un
        \omterm{def:b1:structures:morphism}{morphisme} de \omterm{def:b1:structures:group}{groupes}.
  \item Montrer qu'un $k$-cycle a pour signature $(-1)^{k-1}$, et
        qu'en général $\varepsilon(\sigma) = (-1)^{n -
        c(\sigma)}$, où $c(\sigma)$ est le nombre d'orbites de
        $\sigma$ (points fixes compris).
  \item Le \emph{\omterm{pb:b1:structures:1}{groupe alterné}} est $\mathfrak A_n =
        \ker\varepsilon$. Justifier que c'est un \omterm{def:b1:structures:subgroup}{sous-groupe}, et
        démontrer $\abs{\mathfrak A_n} = \frac{n!}2$ pour $n \geq 2$.
        \emph{(Fixer une transposition $\tau_0$ et considérer $\sigma
        \mapsto \sigma\tau_0$.)}
  \item Vérification de cohérence sur $\sigma = [4, 1, 5, 2, 3, 7, 8, 6]$~:
        calculer $\varepsilon(\sigma)$ de trois façons --- en comptant
        les inversions, à partir du type de cycles via la question 10,
        et à partir du nombre de transpositions de la question 5.
\end{enumerate}

\medskip
\noindent\textbf{Partie III --- $\mathfrak A_n$ est engendré par les
$3$-cycles.}
\begin{enumerate}[resume]
  \item Soient $a, b, c, d$ deux à deux distincts. Vérifier les deux
        identités
        \[
        (a\ b)(a\ c) = (a\ c\ b),
        \qquad
        (a\ b)(c\ d) = (a\ c\ b)(a\ c\ d) .
        \]
  \item Démontrer que pour $n \geq 3$, tout élément de $\mathfrak
        A_n$ est un produit de $3$-cycles. \emph{(Une \omterm{def:b1:counting:objects}{permutation}
        paire est un produit d'un nombre pair de transpositions~;
        les absorber deux par deux.)}
  \item Écrire explicitement $(1\ 2)(3\ 4)$ et le $5$-cycle
        $(1\ 2\ 3\ 4\ 5)$ comme produits de $3$-cycles.
  \item Démontrer la formule de conjugaison~: pour tout $\sigma \in
        \mathfrak S_n$,
        \[
        \sigma\,(a_1\ \dots\ a_k)\,\sigma^{-1}
        = \bigl(\sigma(a_1)\ \dots\ \sigma(a_k)\bigr) .
        \]
\end{enumerate}

\medskip
\noindent\textbf{Partie IV --- Le taquin.} Les pièces $1, \dots, 8$
coulissent dans un cadre $3 \times 3$ comportant une case vide~; un
\emph{coup} fait glisser dans la case vide une pièce qui lui est
adjacente. Numérotons les cases $1, \dots, 9$ (ligne par ligne~; la
position résolue place la pièce $i$ dans la case $i$ et la case vide
en case $9$). Traitons la case vide comme une neuvième pièce~: une
position est alors une \omterm{def:b1:counting:objects}{permutation} $\sigma \in \mathfrak S_9$ (la
pièce $\sigma(i)$ occupe la case $i$).
\begin{enumerate}[resume]
  \item Montrer qu'un coup remplace $\sigma$ par $\sigma \circ \tau$,
        où $\tau$ est la transposition des deux cases concernées~;
        en déduire que chaque coup change le signe de
        $\varepsilon(\sigma)$.
  \item Soit $d(\sigma)$ la distance de Manhattan (lignes plus
        colonnes) entre la case qu'occupe la case vide et sa case
        d'origine $9$. Montrer que chaque coup fait varier $d$ de
        $\pm1$, de sorte que chaque coup change aussi le signe de
        $(-1)^{d(\sigma)}$. En conclure que
        \[
        I(\sigma) = \varepsilon(\sigma)\cdot(-1)^{d(\sigma)}
        \]
        est \emph{invariant} par tout coup.
  \item Démontrer l'impossibilité classique du taquin~: la position
        qui échange les pièces $7$ et $8$ en laissant tout le reste
        (case vide comprise) en place ne peut pas être atteinte à
        partir de la position résolue.
  \item On admet la réciproque (sa démonstration est une récurrence
        instructive mais longue)~: toute position vérifiant $I = +1$
        est atteignable. En déduire qu'exactement la moitié des $8!$
        positions ayant la case vide à sa place sont solubles, soit
        $\frac{8!}2 = 20\,160$.
  \item Déduire de la question 20 que les dispositions de pièces
        atteignables ayant la case vide à sa place forment exactement
        le \omterm{def:b1:structures:subgroup}{sous-groupe} $\mathfrak A_8 \leq \mathfrak S_8$.
  \item Applications de l'invariant~: peut-on atteindre (a) la
        position où les pièces $1, 2, 3$ sont permutées circulairement
        et où tout le reste, case vide comprise, est à sa place~?
        (b) la position où la pièce $5$ et la case vide ont échangé
        leurs places, toutes les autres pièces étant à la leur~?
        Justifier les deux réponses à l'aide de $I$.
\end{enumerate}

\medskip
\noindent\textbf{Partie V --- Synthèse.}
\begin{enumerate}[resume]
  \item Démontrer que pour $n \geq 3$, les seuls \omterm{def:b1:structures:morphism}{morphismes} de
        \omterm{def:b1:structures:group}{groupes} $f \colon \mathfrak S_n \to \{\pm 1\}$ sont le
        \omterm{def:b1:structures:morphism}{morphisme} constant et $\varepsilon$. \emph{(À l'aide de la
        question 16 et de la commutativité de $\{\pm1\}$, montrer que
        $f$ prend la même valeur sur toutes les transpositions.)}
  \item Où exactement le problème a-t-il utilisé~: (i) la notion de
        \omterm{def:b1:structures:morphism}{morphisme} et la \cref{prop:b1:structures:kernel}~; (ii) les
        principes de dénombrement du \cref{ch:b1:counting}~; (iii) le
        problème de bonne définition que les questions 8 et 9
        résolvent~? Une phrase pour chacun.
  \item Synthèse, en un court paragraphe~: une seule fonction de
        parité, dont la bonne définition est démontrée une fois pour
        toutes, organise simultanément la structure interne de
        $\mathfrak S_n$ (le \omterm{def:b1:structures:subgroup}{sous-groupe} $\mathfrak A_n$), tranche un
        casse-tête matériel et --- via la formule $\det = \sum_\sigma
        \varepsilon(\sigma)\cdots$ --- définira les déterminants
        au \cref{ch:b1:det}. Commenter le schéma récurrent~: les
        invariants transforment «~essayer toutes les suites de
        coups~» en un unique calcul.
\end{enumerate}
\end{problem}
